\section{Stein characterization and the marked Laguerre realization}
\label{final:operators}

Throughout this section, $p(t)=te^{-t}$ is the density in the intrinsic
decomposition of $\Sigma$, in the specified real coordinate $t>0$.
The complete identifying data below reconstruct $\Sigma$ with their
coordinate and operator context retained. The candidate classes for a
coordinate probability, an abstract self-adjoint operator, and a
semigroup mixing measure have their respective identifying equations.
All Hilbert spaces are complex.  We use $\mathrm{Id}$ for the identity
operator, and reserve $I$ for the scalar potential.  The arguments give
analytic proofs with explicit domains and candidate classes; numerical
audits are not used to establish these universal statements.

\begin{theorem}[The full-line weak Stein characterization]
\label{final:O1}
Let $P$ be a positive Borel probability measure on $\mathbb R$.  Then
\begin{equation}\label{final:O1-identity}
 \int_{\mathbb R}\bigl[t f'(t)+(2-t)f(t)\bigr]P(dt)=0
 \qquad\bigl(f\in C_c^\infty(\mathbb R)\bigr)
\end{equation}
if and only if
\[
 P(dt)=\mathbf 1_{(0,\infty)}(t)te^{-t}\,dt.
\]
In particular, absolute continuity, positive support, absence of an atom
at zero, and all finite moments are conclusions, not hypotheses.
\end{theorem}

\begin{proof}
The integrands in \eqref{final:O1-identity} have compact support, so no
moment assumption is needed to state the identity.  In distributions it
says $(tP)'=(2-t)P$.  On either open half-line put $Q=tP$.  There
\[
 Q'=(2/t-1)Q,\qquad
 \bigl(e^t t^{-2}Q\bigr)'=0.
\]
A distribution with zero derivative on a connected interval is a constant
distribution.  Thus the restrictions of $P$ have densities
$C_+te^{-t}$ on the positive half-line and $C_-|t|e^{-t}$ on the negative
half-line, with $C_+,C_-\geq0$.  Finiteness of $P$ forces $C_-=0$.
The only remaining possible part is $m\delta_0$, $m\geq0$.
The positive-half-line flux $t^2e^{-t}$ tends to zero at the origin;
its distributional derivative therefore has no atom there.  Substitution
of $m\delta_0$ into the equation gives $0=2m\delta_0$, so $m=0$.
Finally $\int_0^\infty te^{-t}\,dt=1$ gives $C_+=1$.
Conversely, integration by parts, using
$(t^2e^{-t})'=(2-t)te^{-t}$ and the zero endpoint flux, proves the
displayed identity.  Tests supported only in $(0,\infty)$ could not
exclude mass on $(-\infty,0]$; the full-line test class is essential
to this support conclusion.
\end{proof}

\begin{proposition}[The integrable centered Stein kernel]
\label{final:O1-kernel}
Fix the probability density $p$, and let $\tau$ be real measurable with
$\tau p\in L^1_{\mathrm{loc}}(0,\infty)$.  The equation
\[
 (\tau p)'=(2-t)p
\]
in distributions is equivalent to
\begin{equation}\label{final:O1-kernel-family}
 \tau(t)=t+C\frac{e^t}{t}\quad\hbox{almost everywhere},
 \qquad C\in\mathbb R.
\end{equation}
The locally absolutely continuous representative of $\tau p$ has both
endpoint limits equal to $C$.  Each of the following separately selects
$\tau=t$ almost everywhere: zero flux at either endpoint;
$\tau\in L^1(p\,dt)$; a finite expectation of $\tau$ under $p\,dt$; or
the finite identity $\int\tau p=\operatorname{Var}_p(t)=2$.
Nonnegativity alone permits exactly $C\geq0$.
\end{proposition}

\begin{proof}
Write $q=tp=t^2e^{-t}$.  Since $q'=(2-t)p$, the distribution
$\tau p-q$ has derivative zero and is constant.  This proves
\eqref{final:O1-kernel-family} and the assertions about its absolutely
continuous representative.  For $C\ne0$, $q+C$ tends to $C$ at infinity,
so $\int_0^\infty|\tau|p=\int_0^\infty|q+C|=\infty$; its tail also
has the constant sign of $C$, so $\int \tau\,p\,dt$ cannot be finite as
a Lebesgue integral.  For $C=0$, integration by parts gives
$\int tp=2$, $\int t^2p=6$, and hence variance $2$.
As $q>0$ and $q$ tends to zero at both endpoints, $q+C\geq0$
almost everywhere exactly when $C\geq0$.

Every value of $C$ satisfies the compact-test identity
\[
 \int\tau f'p\,dt=\int(t-2)fp\,dt
 \qquad(f\in C_c^\infty(0,\infty)),
\]
because the added term is $C\int f'=0$.  Including $f(t)=t$, with
both integrals required finite, instead gives $\int\tau p=6-4=2$.
This inverse problem fixes $P$ and varies $\tau$; it is different from
Theorem~\ref{final:O1}, which fixes the coefficient $t$ and varies $P$.
\end{proof}

\begin{theorem}[Two marked probes and the Pearson realization]
\label{final:O2}
In the conservative local second-order class
\[
 Bf=-a(t)f''-b(t)f',\qquad t>0,
\]
the two differential-expression identities
\begin{equation}\label{final:O2-probes}
 B(2-t)=2-t,\qquad
 B(t^2/2-3t+3)=t^2-6t+6
\end{equation}
hold if and only if $a(t)=t$ and $b(t)=2-t$.
The conclusion is pointwise or almost everywhere according to the
meaning of the coefficient identities.  If a positive, locally
absolutely continuous probability density $w$ is also required to obey
the zero-current Pearson equation $(aw)'=bw$, then $w=p$.
Thus these data determine the differential expression and its normalized
Hilbert weight in this specified coordinate class.
\end{theorem}

\begin{proof}
The first probe gives $b=2-t$.  The second then gives
\[
 -a-(2-t)(t-3)=t^2-6t+6,
\]
and consequently $a=t$.  Substitution proves the converse.
The Pearson equation reduces to
$tw'=(1-t)w$, so $(e^t w/t)'=0$ locally and $w=Cte^{-t}$.
Probability normalization gives $C=1$.
The probes are polynomials; they are identities of the expression,
or identities on an explicitly enlarged domain containing those
polynomials.  They are not assertions that polynomials belong to
$C_c^\infty(0,\infty)$.
\end{proof}

\begin{proposition}[The exact limits of the probe and stationary data]
\label{final:O2-boundaries}
The pair \eqref{final:O2-probes} is irredundant within the indicated
coefficient class, even if the coefficients are smooth and $a>0$.
It does not identify an arbitrary self-adjoint operator outside that
local class.  The invariant density $p$ alone does not identify a
reversible dynamics or its time scale.
\end{proposition}

\begin{proof}
With only the linear probe, set $b=2-t$ and choose any positive smooth
$a$.  With only the quadratic probe, take a nonzero
$k\in C_c^\infty(0,\infty)$ and set
\[
 b=2-t+\varepsilon k(t),\qquad
 a=t-\varepsilon k(t)(t-3).
\]
For sufficiently small $|\varepsilon|$, positivity holds on the compact
support of $k$ and elsewhere $a=t$; the quadratic identity holds
identically while the linear one changes.
For the larger self-adjoint category, use the operator and eigenbasis
constructed below.  Let $U$ fix $e_0,e_1,e_2$ and every $e_n$, $n\geq5$,
and make a nontrivial rotation, with angle in $(0,\pi/2)$, on
$\operatorname{span}\{e_3,e_4\}$.  Then $UAU^{-1}\ne A$, since the
two rotated eigenvalues are distinct, but it preserves the constant
vector and both probes.  This example has self-adjointness; positivity
preservation of the rotated semigroup is not imposed.
Finally $cA$, $c>0$, gives distinct conservative reversible semigroups
with the same invariant $p\,dt$, as follows from the construction below.
The centered integrable Stein selection fixes $c=1$; invariance alone
does not.
\end{proof}

\subsection{The actual self-adjoint domain and complete spectral resolution}

Set
\[
 \mathcal H=L^2((0,\infty),p(t)\,dt),\qquad q(t)=tp(t)=t^2e^{-t},
 \qquad \langle f,g\rangle=\int_0^\infty f\overline g\,p\,dt.
\]
In particular $\langle\cdot,\cdot\rangle$ is linear in its first
argument.  The minimal differential operator considered here is
\begin{equation}\label{final:O3-minimal}
 A_0f=-p^{-1}(qf')'=-tf''-(2-t)f',\qquad
 D(A_0)=C_c^\infty(0,\infty).
\end{equation}

\begin{theorem}[Essential self-adjointness and the full domain]
\label{final:O3}
The operator $A_0$ is densely defined, nonnegative and symmetric.
Both singular endpoints are limit point, its deficiency indices are
$(0,0)$, and its unique self-adjoint extension is
$A=\overline{A_0}=A_0^*$.  Its exact domain is
\begin{equation}\label{final:O3-domain}
 D(A)=\left\{f\in\mathcal H:
 f,qf'\in AC_{\mathrm{loc}}(0,\infty),\quad
 -p^{-1}(qf')'\in\mathcal H\right\}.
\end{equation}
Here the regularity refers to the indicated representatives of the
$L^2$ classes.  For every $f\in D(A)$,
\begin{equation}\label{final:O3-flux}
 \lim_{t\downarrow0}q(t)f'(t)=
 \lim_{t\to\infty}q(t)f'(t)=0.
\end{equation}
These limits are consequences of the domain, not additional boundary
conditions.  In particular no finite boundary value $f(0)$ is required.
\end{theorem}

\begin{proof}
Smooth compactly supported functions are dense in this weighted $L^2$
space, because the weight is positive and smooth on each compact
interior interval.  Integration by parts gives
\[
 \langle A_0f,g\rangle=\int_0^\infty qf'\overline{g'}\,dt,
\]
which proves symmetry and nonnegativity.
The definition of the adjoint says that $f\in D(A_0^*)$ precisely when
there is an $h\in\mathcal H$ with $-(qf')'=ph$ in distributions.
The right side is locally integrable.  Taking a primitive, dividing
by the positive smooth function $q$, and taking a second primitive
shows that $qf'$ and $f$ have locally absolutely continuous
representatives.  Conversely those representatives permit integration
by parts against each compact test.  Thus the adjoint is exactly the
maximal operator in \eqref{final:O3-domain}.

To determine the endpoints, two independent zero-energy solutions are
\[
 u_1(t)=1,\qquad u_2(t)=\int_1^t\frac{e^s}{s^2}\,ds.
\]
Expanding $e^s/s^2=s^{-2}+s^{-1}+O(1)$ near zero and using
l'H\^opital's rule at infinity gives
\[
 u_2(t)=-t^{-1}+\log t+O(1)\quad(t\downarrow0),\qquad
 u_2(t)\sim e^t/t^2\quad(t\to\infty).
\]
Consequently $|u_2|^2p\sim t^{-1}$ at zero and
$|u_2|^2p\sim e^t/t^3$ at infinity.  Neither is integrable.
The constant solution is square integrable near each endpoint.
The Weyl endpoint alternative therefore makes both endpoints limit
point: an endpoint is limit circle precisely when all solutions at
one, equivalently every, spectral parameter are square integrable
there.  The deficiency-index theorem gives one deficiency at each
limit-circle endpoint and zero at each limit-point endpoint.
Its hypotheses hold here because $p,q>0$, $p,q^{-1}$ are locally
integrable, and the potential is zero; see
Gesztesy--Littlejohn--Nichols~\cite[Theorems 3.6 and 3.8]{GLN}.
Their minimal domain is the closure of compactly
supported maximal-domain elements.  It has the same closure as
$C_c^\infty$ here: on compact interior intervals those elements are
in $H^2$, and smooth approximation converges in the graph norm.
This verifies the application to the actual domain
\eqref{final:O3-minimal}, including the critical logarithmic divergence
at the origin.  It gives $\overline{A_0}=A_0^*$ and deficiency indices
$(0,0)$.  Nonnegativity passes to this closure.

For the flux statement put $h=Af$.  Cauchy--Schwarz and
$\int p=1$ give $\int|h|p\leq\|h\|_{\mathcal H}$.
Thus $(qf')'=-ph$ is globally integrable and $qf'$ has finite limits
at both endpoints.  A nonzero limit $\ell$ at zero would yield
$f(t)\sim-\ell/t$, by integrating $f'=(\ell+o(1))/q$;
a nonzero limit at infinity would give
$f(t)\sim\ell e^t/t^2$.  The same weighted-square divergences exclude
both.  This proves \eqref{final:O3-flux}; in fact
\[
 q(t)f'(t)=-\int_0^t hp\,ds=\int_t^\infty hp\,ds,
 \qquad \int_0^\infty hp\,dt=0.
\]
To see why finiteness of $f(0)$ cannot be added, multiply
$\log\log(1/t)$ by a smooth cutoff supported sufficiently near zero
and equal to one there.  This function is in $\mathcal H$, and near
zero its image under the differential expression is
$O((t\log(1/t))^{-1})$.  Its weighted square integral is finite,
since $\int_0^\varepsilon dt/(t\log^2(1/t))<\infty$.
It therefore belongs to \eqref{final:O3-domain} and has no finite
value at zero.
\end{proof}

\begin{proposition}[The closed form and the conservative semigroup]
\label{final:O3-form}
The closed form of $A$ is
\[
 \mathcal E(f,g)=\int_0^\infty qf'\overline{g'}\,dt,
\]
with the exact domain
\begin{equation}\label{final:O3-form-domain}
 D(A^{1/2})=
 \left\{f\in\mathcal H:f\in AC_{\mathrm{loc}}(0,\infty),\quad
 \int_0^\infty q|f'|^2\,dt<\infty\right\}.
\end{equation}
There are no additional endpoint traces.  The semigroup $e^{-sA}$,
$s\geq0$, is positivity preserving, preserves the constant function
$1$, and preserves the probability $p\,dt$.
\end{proposition}

\begin{proof}
Denote the displayed form domain by $V$.  A Cauchy sequence in
$\|f\|_V^2=\|f\|_{\mathcal H}^2+\int q|f'|^2$ is Cauchy in
ordinary $H^1$ on every compact interior interval, where both weights
are bounded above and below by positive constants.  Its local weak
derivative equals the weighted derivative limit, proving completeness.
Truncate real and imaginary parts of $f\in V$ at height $N$.
The weak chain rule and dominated convergence show convergence in
$V$ as $N\to\infty$.  For bounded $f$, choose cutoffs which rise on
$[\varepsilon,2\varepsilon]$, fall on $[R,2R]$, and have derivatives
bounded respectively by $C/\varepsilon$ and $C/R$.  Their energy is
bounded by
\[
 C\varepsilon+\frac{C}{R^2}\int_R^{2R}t^2e^{-t}\,dt\longrightarrow0.
\]
The two derivative errors in $\chi f-f$ are
$(\chi-1)f'$ and $\chi'f$; the first tends to zero by dominated
convergence, the second by this estimate and boundedness of $f$.
The $\mathcal H$ norm also tends to zero.  On the resulting compact
interior support, mollification converges in $H^1$ and hence in $V$.
Thus $C_c^\infty$ is a form core, and the displayed form is the
closure of the minimal form.  Its represented self-adjoint operator
extends $A_0$ and must equal $A$ by Theorem~\ref{final:O3}.

The weak chain rule shows that real normal contractions preserve $V$
and decrease its energy.  Equivalently, the variational resolvent is
positivity preserving and maps $[0,1]$ into $[0,1]$, as follows by
testing its equation with the negative part and with $(u-1)^+$.
The exponential resolvent limit gives the same Markov property for
$e^{-sA}$.  Finally $1\in D(A)$ and $A1=0$, so $e^{-sA}1=1$.
Self-adjointness gives
$\int e^{-sA}f\,p\,dt=\langle e^{-sA}f,1\rangle
=\langle f,e^{-sA}1\rangle=\int fp\,dt$.
\end{proof}

\begin{theorem}[The Laguerre basis and the exact integer spectrum]
\label{final:O3-spectrum}
Define the polynomials with their marked normalization by
\begin{equation}\label{final:O3-laguerre}
 L_n^{(1)}(t)=\frac{e^t}{n!t}\frac{d^n}{dt^n}
       (e^{-t}t^{n+1})
 =\sum_{k=0}^n(-1)^k\binom{n+1}{k+1}\frac{t^k}{k!},
 \qquad n\in\mathbb N_0.
\end{equation}
Then $e_n=L_n^{(1)}/\sqrt{n+1}$ is a complete orthonormal basis of
$\mathcal H$, each $e_n\in D(A)$, and $Ae_n=ne_n$.
In particular the spectrum is exactly $\mathbb N_0$, each eigenvalue
is simple, and the resolvent is compact.  If $f=\sum c_ne_n$, then
\begin{align}
 D(A)&=\left\{f:\sum_{n=0}^\infty n^2|c_n|^2<\infty\right\},
 &Af&=\sum_{n=0}^\infty nc_ne_n,\label{final:O3-spectral-domain}\\
 D(A^{1/2})&=\left\{f:\sum_{n=0}^\infty n|c_n|^2<\infty\right\},
 &\mathcal E(f,f)&=\sum_{n=0}^\infty n|c_n|^2.
 \nonumber
\end{align}
\end{theorem}

\begin{proof}
Leibniz's formula gives the finite sum in
\eqref{final:O3-laguerre}.  Its coefficients $a_k$ satisfy
$(k+1)(k+2)a_{k+1}+(n-k)a_k=0$, so
\[
 t(L_n^{(1)})''+(2-t)(L_n^{(1)})'+nL_n^{(1)}=0.
\]
Every polynomial belongs to the maximal domain: the polynomial and
its polynomial image are in $\mathcal H$, and all local regularity
conditions hold.  These are therefore genuine eigenfunctions of $A$.
For $m<n$, integrating the Rodrigues formula $n$ times by parts gives
$\int pL_n^{(1)}t^m\,dt=0$.  The boundary terms vanish at infinity
by exponential decay and at zero because derivatives of
$e^{-t}t^{n+1}$ of order at most $n-1$ vanish to order at least two.
For $m=n$ the integral is $(-1)^n\Gamma(n+2)$.
The leading coefficient $(-1)^n/n!$ then gives
\[
 \int_0^\infty L_n^{(1)}L_m^{(1)}p\,dt=(n+1)\delta_{nm}.
\]

For completeness, let $g\in\mathcal H$ be orthogonal to every
polynomial and define
$F(z)=\int_0^\infty g(t)p(t)e^{zt}\,dt$.
Cauchy--Schwarz gives locally uniform domination of this integral
and of all its differentiated integrands on $\operatorname{Re}z<1/2$.
Thus $F$ is holomorphic there, and $F^{(n)}(0)=0$ for all $n$.
The identity theorem gives $F=0$ on that half-plane.
The finite complex measure $gp\,dt$, pushed forward by $y=e^{-t}$
and regarded as a measure on $[0,1]$, consequently has every moment
zero, since $F(-k)=0$ for $k\in\mathbb N_0$.
Uniform polynomial approximation on $[0,1]$ shows that it integrates
every continuous function to zero; measure uniqueness gives $g=0$.
This proves completeness.
The unitary coefficient map now takes $A$ to multiplication by $n$
on $\ell^2(\mathbb N_0)$, with exactly the displayed domain.
The form domain follows by taking square roots.  The diagonal
entries $(n+\alpha)^{-1}$ of any resolvent, $\alpha>0$, tend to zero,
proving compactness and the asserted spectrum.
\end{proof}

\begin{theorem}[Complete marked Laguerre orthogonality identifies the law]
\label{final:O7}
Let $P$ be a positive Borel probability measure on $\mathbb R$ for
which every polynomial is integrable.  With the specific polynomials
\eqref{final:O3-laguerre},
\[
 \int_{\mathbb R}L_n^{(1)}(t)\,P(dt)=0\quad(n\geq1)
\]
holds if and only if $P(dt)=\mathbf1_{(0,\infty)}p(t)\,dt$.
The positive support is again a conclusion.
\end{theorem}

\begin{proof}
Mass fixes moment zero.  The nonzero leading coefficient
$(-1)^n/n!$ means that the $n$th displayed equation determines moment
$n$ from all smaller moments.  The preceding Rodrigues calculation
proves these equations for $p\,dt$, whose moments are $(n+1)!$.
Induction therefore gives $\int t^nP(dt)=(n+1)!$ for all $n\geq0$.
For $0<c<1$, monotone convergence in the nonnegative even-power
series gives
\[
 \int\cosh(ct)\,P(dt)
 =\sum_{n=0}^\infty(2n+1)c^{2n}
 =\frac{1+c^2}{(1-c^2)^2}<\infty.
\]
As $e^{c|t|}\leq2\cosh(ct)$, the bilateral moment-generating
function is holomorphic on $|\operatorname{Re}z|<1$.
Its Taylor series near zero, determined by the moments, is
$\sum_{n\geq0}(n+1)z^n=(1-z)^{-2}$.
The identity theorem gives equality to the Gamma moment-generating
function throughout the strip, including the imaginary axis.
Uniqueness of characteristic functions identifies $P$ on all of
$\mathbb R$.  The converse was already proved by Rodrigues's formula.
\end{proof}

\subsection{Complete spectral data, kernel formulas, and their inverse category}

\begin{theorem}[Exact trace domains and reconstruction of the unitary class]
\label{final:O4}
For the operator $A$ above, the ordinary traces are
\begin{align}
 Z_A(s)&=\operatorname{Tr}(\mathrm{Id}+A)^{-s}=\zeta(s),
 &&\operatorname{Re}s>1,\label{final:O4-zeta}\\
 \Theta_A(\tau)&=\operatorname{Tr}e^{-\tau A}
       =\frac1{1-e^{-\tau}},
 &&\operatorname{Re}\tau>0.\label{final:O4-heat-trace}
\end{align}
These are precisely their trace-class domains.  The powers use the
real logarithms of the positive eigenvalues $n+1$; the zero eigenspace
is included through this shift.

Let $B\geq0$ be self-adjoint with compact resolvent.  Each of the
following complete data sets determines its eigenvalue multiset,
including multiplicities, and hence its unitary equivalence class:
all finite heat traces $\operatorname{Tr}e^{-\tau B}$, $\tau>0$;
all shifted zeta traces on a real half-line $s>s_0$ on which they are
finite; or all integer shifted traces
$\operatorname{Tr}(\mathrm{Id}+B)^{-k}$, $k\geq2$, with
$\operatorname{Tr}(\mathrm{Id}+B)^{-2}<\infty$.
The coordinate law of an unspecified realization is not part of this
unitary-class conclusion.
\end{theorem}

\begin{proof}
The basis of Theorem~\ref{final:O3-spectrum} gives diagonal entries
$(n+1)^{-s}$ and $e^{-\tau n}$.  Their singular values are respectively
$(n+1)^{-\operatorname{Re}s}$ and $e^{-n\operatorname{Re}\tau}$.
Their sums converge exactly on the stated half-planes and there give
the Dirichlet and geometric series in
\eqref{final:O4-zeta}--\eqref{final:O4-heat-trace}.
Analytic or meromorphic continuations outside these domains do not
become ordinary operator traces.

For the inverse, compact resolvent gives an orthonormal eigenbasis,
finite multiplicities, and eigenvalues tending to infinity if there
are infinitely many.  For any nonempty remaining multiset let its
least eigenvalue be $\lambda_0$, of multiplicity $m_0$.
If $T(\tau)$ is its remaining heat trace, then
\[
 \lambda_0=-\lim_{\tau\to\infty}\frac{\log T(\tau)}{\tau},
 \qquad m_0=\lim_{\tau\to\infty}e^{\tau\lambda_0}T(\tau).
\]
Indeed, after this rescaling the least-eigenvalue terms equal $m_0$,
and the other terms tend to zero.  Finiteness at a fixed positive
time supplies a summable dominating sequence for all larger times.
Subtract $m_0e^{-\tau\lambda_0}$ and iterate.  The identically zero
remainder detects the end of a finite spectrum, including the empty
Hilbert space.
For zeta data write $r_j=1+\lambda_j$.  The analogous formulas are
\[
 r_0=\lim_{s\to\infty}Z(s)^{-1/s},\qquad
 m_0=\lim_{s\to\infty}r_0^sZ(s).
\]
Finiteness at one exponent again dominates the rescaled tails.
Subtracting $m_0r_0^{-s}$ and repeating recovers every eigenvalue.

For integer data set $x_j=(1+\lambda_j)^{-1}$ and form the finite
positive measure
\[
 \eta_B=\sum_j x_j^2\delta_{x_j}\quad\hbox{on }[0,1],
 \qquad \int x^n\eta_B(dx)=Z_B(n+2),\quad n\geq0.
\]
Two such measures with the same moments agree on polynomials and
then on all continuous functions by uniform polynomial approximation.
They are therefore equal.  Every nonzero atom at $x$ has mass
$m x^2$, recovering the integer multiplicity $m$ and eigenvalue
$x^{-1}-1$.  There is no atom at zero from any finite eigenvalue.
This is the compact Hausdorff moment argument; it uses no interpolation
from the integer values.  Finally an eigenvalue-preserving bijection
of orthonormal eigenbases is an intertwining unitary, with the
weighted coefficient domains carried to one another.
\end{proof}

\begin{proposition}[Compact resolvent does not supply finite trace data]
\label{final:O4-convergence-boundary}
There is a nonnegative diagonal self-adjoint operator with compact
resolvent for which neither a heat trace nor a positive shifted zeta
trace is finite at any positive parameter.
\end{proposition}

\begin{proof}
On $\ell^2(\mathbb N)$ take eigenvalues
$\lambda_j=\log\log(j+e^e)$, with domain
$\{c:\sum\lambda_j^2|c_j|^2<\infty\}$.
They tend to infinity, so $(\mathrm{Id}+B)^{-1}$ is compact.
For every $\tau>0$, $(\log(j+e^e))^{-\tau}$ eventually exceeds
$1/j$, and for every $s>0$,
$(1+\log\log(j+e^e))^{-s}$ eventually exceeds $1/j$.
Both trace series diverge.  Thus the convergence hypotheses in
Theorem~\ref{final:O4} cannot be inferred from compactness alone.
\end{proof}

\begin{theorem}[The complete heat kernel relative to the invariant measure]
\label{final:O4-heat-kernel}
For $\tau>0$, set $r=e^{-\tau}$.  The kernel of $e^{-\tau A}$ relative
to $p(y)\,dy$ is
\begin{align}
 K_\tau(x,y)
 &=\sum_{n=0}^\infty\frac{r^n}{n+1}
       L_n^{(1)}(x)L_n^{(1)}(y)\label{final:O4-kernel-series}\\
 &=\frac{\exp[-r(x+y)/(1-r)]}{(1-r)\sqrt{rxy}}
       I_1\!\left(\frac{2\sqrt{rxy}}{1-r}\right),
       \qquad x,y>0.\label{final:O4-hille-hardy}
\end{align}
Here $I_1$ is the modified Bessel function, specified by
$I_1(z)=\sum_{m\geq0}(z/2)^{2m+1}/(m!(m+1)!)$.
The series converges locally uniformly in $x,y>0$ and in
$L^2(p(x)p(y)\,dx\,dy)$.
The kernel is positive and symmetric,
$\int K_\tau(x,y)p(y)\,dy=1$, and it tends to $1$ as
$\tau\to\infty$.  The Lebesgue transition density is
$K_\tau(x,y)p(y)$.
\end{theorem}

\begin{proof}
Here is a direct verification of the normalization and the closed
formula, which is the parameter-one Hille--Hardy
identity~\cite[\href{https://dlmf.nist.gov/18.18.E27}{Eq.~18.18.27}]{DLMF}.
The finite coefficients in \eqref{final:O3-laguerre} give the
generating function
\[
 G_z(y)=\sum_{n\geq0}L_n^{(1)}(y)z^n
       =(1-z)^{-2}\exp[-yz/(1-z)],\qquad |z|<1.
\]
Let $\widetilde K$ be the right side of
\eqref{final:O4-hille-hardy}.  Expanding $I_1$ gives the nonnegative
series
\[
 \widetilde K(x,y)=e^{-r(x+y)/(1-r)}
   \sum_{m=0}^\infty
   \frac{(rxy)^m}{m!(m+1)!(1-r)^{2m+2}}.
\]
For real $z$ sufficiently near zero, put
$c=(1-rz)/((1-r)(1-z))>0$.
Termwise gamma integration, justified by positivity, gives
\begin{align*}
 \int_0^\infty\widetilde K(x,y)G_z(y)p(y)\,dy
 &=\frac{e^{-rx/(1-r)}}{(1-z)^2(1-r)^2c^2}
       \exp\!\left(\frac{rx}{(1-r)^2c}\right)\\
 &=(1-rz)^{-2}\exp[-rxz/(1-rz)]=G_{rz}(x).
\end{align*}
In particular $z=0$ gives row mass one.  The kernel is symmetric and
positive, so Jensen's inequality and row mass one make its integral
operator a contraction on $\mathcal H$.  For $z$ in a sufficiently
small complex neighborhood of zero, the same exponentially dominated
integral is holomorphic.  Differentiating at zero proves that its
operator maps $L_n^{(1)}$ to $r^nL_n^{(1)}$.  Completeness of the
basis identifies it with $e^{-\tau A}$.
The spectral series is its $L^2$ kernel, since
$\sum r^{2n}<\infty$.  Cauchy's coefficient estimate for $G_z$,
on any circle $|z|=R$ with $\sqrt r<R<1$, bounds its summands
uniformly on compact sets by a convergent geometric sequence.
Thus this series is continuous and locally uniformly convergent.
Its $L^2$ equality with the continuous closed formula implies their
pointwise equality on $(0,\infty)^2$.  Finally the $m=0$ Bessel term
and the convergent series give $K_\tau(x,y)\to1$ as $r\downarrow0$.
The definition of an integral kernel relative to $p(y)\,dy$ proves
the assertion about the Lebesgue density.
\end{proof}

\begin{proposition}[Resolvents, their kernels, and the operator inverse]
\label{final:O4-resolvents}
For $\alpha>0$, the kernels relative to $p(y)\,dy$ of
$(A+\alpha\mathrm{Id})^{-1}$ and $(A+\alpha\mathrm{Id})^{-2}$ are
respectively
\begin{align*}
 R_\alpha(x,y)&=\int_0^\infty e^{-\alpha\tau}K_\tau(x,y)\,d\tau,\\
 S_\alpha(x,y)&=\int_0^\infty \tau e^{-\alpha\tau}K_\tau(x,y)\,d\tau.
\end{align*}
These identities hold as $L^2$ kernels, and the nonnegative integrals
are finite almost everywhere.  The first resolvent is Hilbert--Schmidt
but is not trace class.  The second is trace class and
\[
 \operatorname{Tr}(A+\alpha\mathrm{Id})^{-2}=\zeta(2,\alpha),
\]
where $\zeta(s,\alpha)=\sum_{n=0}^\infty(n+\alpha)^{-s}$ for
$\operatorname{Re}s>1$.  For $\alpha,\beta>0$,
\begin{equation}\label{final:O4-resolvent-difference}
 \operatorname{Tr}\bigl[(A+\alpha\mathrm{Id})^{-1}
       -(A+\beta\mathrm{Id})^{-1}\bigr]
       =\psi(\beta)-\psi(\alpha),
\end{equation}
with $\psi=\Gamma'/\Gamma$.
A complete resolvent operator, with its Hilbert realization and shift
retained, determines $A$ and its domain.  A full marked heat operator
at any fixed positive time also determines $A$.
\end{proposition}

\begin{proof}
On each eigenvector the Laplace integrals are
$\int_0^\infty e^{-(n+\alpha)\tau}\,d\tau=(n+\alpha)^{-1}$ and
$\int_0^\infty\tau e^{-(n+\alpha)\tau}\,d\tau=(n+\alpha)^{-2}$.
The corresponding strong operator integrals converge, and their
kernel integrals converge in $L^2$ because
\[
 \|K_\tau\|_{L^2(p\otimes p)}
       =(1-e^{-2\tau})^{-1/2},
\]
which is $O(\tau^{-1/2})$ near zero and bounded at infinity.
The exponential weights are therefore integrable in this norm.
Positivity and Tonelli identify these $L^2$ limits with the
nonnegative pointwise integrals almost everywhere.
The singular-value series for the resolvents are
$\sum(n+\alpha)^{-1}=\infty$ and
$\sum(n+\alpha)^{-2}<\infty$, proving the ideal statements.
The difference series is absolutely summable and equals
\[
 \sum_{n=0}^\infty
   \left(\frac1{n+\alpha}-\frac1{n+\beta}\right)
 =\psi(\beta)-\psi(\alpha),
\]
as follows by subtracting the convergent digamma expansions
$\psi(a)=-\gamma_E+\sum_{n\geq0}((n+1)^{-1}-(n+a)^{-1})$.
This also fixes the sign in \eqref{final:O4-resolvent-difference}.

For $R=(A+\alpha\mathrm{Id})^{-1}$ one has
$D(A)=\operatorname{Ran}R$ and
$A(Rh)=h-\alpha Rh$ for every $h\in\mathcal H$.
Thus the complete operator recovers the generator, whereas a scalar
resolvent trace discards its eigenvectors and coordinate realization.
For $T=e^{-\tau A}$ the Borel functional calculus similarly gives
$A=-\tau^{-1}\log T$, with its spectral domain; zero is not an
eigenvalue of $T$, and its possible accumulation there is handled by
this unbounded functional calculus.  Consequently a full kernel with
its reference measure retained recovers the corresponding operator.
\end{proof}

\begin{proposition}[Zeta and Fredholm determinants with their conventions]
\label{final:O4-determinants}
For $\alpha>0$, define the zeta determinant using the meromorphic
continuation of the Hurwitz zeta function.  Then
\begin{equation}\label{final:O4-zeta-determinant}
 \det{}_{\zeta}(A+\alpha\mathrm{Id})
 =\exp[-\partial_s\zeta(0,\alpha)]
 =\frac{\sqrt{2\pi}}{\Gamma(\alpha)}.
\end{equation}
In particular
$\det{}_{\zeta}(\mathrm{Id}+A)=\det{}'_{\zeta}(A)=\sqrt{2\pi}$,
where the prime omits the zero mode.
For every $z\in\mathbb C$ one has the entire-function identities
\begin{align}
 \det\bigl(\mathrm{Id}+z(\mathrm{Id}+A)^{-2}\bigr)
 &=\prod_{k=1}^\infty(1+z/k^2)
   =\frac{\sinh(\pi\sqrt z)}{\pi\sqrt z},
   \label{final:O4-fredholm}\\
 \det{}_2\bigl(\mathrm{Id}+z(\mathrm{Id}+A)^{-1}\bigr)
 &=\prod_{k=1}^\infty(1+z/k)e^{-z/k}
   =\frac{e^{-\gamma_E z}}{\Gamma(1+z)}.
   \label{final:O4-det2}
\end{align}
The first determinant is the ordinary trace-class Fredholm determinant;
the second is the Hilbert--Schmidt regularized determinant.
The square-root quotient denotes its entire power series, with
value $1$ at zero.  Its value is independent of the square-root choice.
The full zero multiset of either displayed determinant identifies
the eigenvalue multiset of $A$ in the corresponding nonnegative,
shifted determinant category.
\end{proposition}

\begin{proof}
The spectral series for $A+\alpha\mathrm{Id}$ is the Hurwitz zeta
series on $\operatorname{Re}s>1$.  Its continuation is regular at
zero and obeys
$\partial_s\zeta(0,\alpha)=\log\Gamma(\alpha)-\tfrac12\log(2\pi)$;
see \cite[\href{https://dlmf.nist.gov/25.11.E18}{Eq.~25.11.18}]{DLMF}.
This gives \eqref{final:O4-zeta-determinant}.  The positive eigenvalues
of $A$ and all eigenvalues of $\mathrm{Id}+A$ are both $1,2,\ldots$,
which proves the two special determinant values with the stated
zero-mode convention.
The eigenvalues of $(\mathrm{Id}+A)^{-2}$ are $k^{-2}$, a summable
sequence; those of $(\mathrm{Id}+A)^{-1}$ are $k^{-1}$, whose squares
are summable.  The determinant definitions therefore give the two
products, locally uniformly in $z$.  The identities
\cite[\href{https://dlmf.nist.gov/4.36.E1}{Eqs.~4.36.1 and
5.8.2}]{DLMF} evaluate them as shown.
Evenness of $\sinh w/w$ proves the branch statement.
For a nonnegative operator $B$ with the respective Schatten-class
shifted inverse, a nonzero eigenvalue $\kappa$ of that inverse
produces a zero at $-1/\kappa$ with its multiplicity.
The maps $\lambda\mapsto(1+\lambda)^{-2}$ and
$\lambda\mapsto(1+\lambda)^{-1}$ are injective on $[0,\infty)$.
The shift, exponent, and regularization convention therefore give
the claimed inverse from the entire zero multiset.
\end{proof}

\begin{theorem}[Finite familiar spectral invariants do not determine the spectrum]
\label{final:O4-finite-countermodels}
Fix any finite list of values of the continued shifted zeta function
at nonzero real regular points, optionally its finite part at $s=1$
and its zeta determinant.  There is a nonconstant local family of
nonnegative self-adjoint diagonal operators, differing from $A$ in
only finitely many positive eigenvalues, that retains all these
specified invariants.  Their zeta functions also retain every pole
and principal part, and the value at zero.  Thus this finite list,
including a single zeta determinant as a special case, does not
identify the complete spectrum.
\end{theorem}

\begin{proof}
Vary $N$ distinct shifted eigenvalues $r_j=1+\lambda_j>1$ and fix
the others.  The correction
\[
 \Delta Z(s)=\sum_{j=1}^N(r_j^{-s}-r_{j,0}^{-s})
\]
is entire, leaves every principal part unchanged, and has
$\Delta Z(0)=0$.  Preserving a value at a nonzero real $s_i$ imposes
the equation $\sum r_j^{-s_i}=\sum r_{j,0}^{-s_i}$.
Preserving the finite part at one imposes this same equation with
$s_i=1$.  Preserving the determinant imposes
$\sum\log r_j=\sum\log r_{j,0}$.
After duplicate equations are removed, let their number be $m$.
The gradient rows are nonzero multiples of $r_j^{-s_i-1}$,
and, if present, the additional row $r_j^{-1}$.
Their exponents are distinct.

These rows have rank $m$ at distinct positive $r_j$ whenever $N\geq m$.
Indeed, put $u_j=\log r_j$.  A nonzero linear combination of $m$
exponentials with distinct real exponents has at most $m-1$ distinct
real zeros.  Inductively, divide by its smallest exponential and
differentiate; Rolle's theorem would otherwise give at least $m-1$
zeros for a nonzero combination of $m-1$ exponentials.  The
single-exponential case starts the induction.  Thus a nonzero row
combination cannot vanish at $m$ distinct $u_j$.
Choose $N>m$.  The implicit-function theorem gives a local level
manifold of dimension $N-m>0$ through the original shifted
eigenvalues.  In small disjoint neighborhoods they remain greater
than one and distinct; changes in this manifold cannot merely
permute the spectrum.  The fixed infinite tail preserves compact
resolvent and every trace convergence domain.

For an explicit simultaneous example replace $2,3,4$ by $a,b,c$,
where
\[
 a=2+\varepsilon,\qquad b+c=7-\varepsilon,
 \qquad bc=\frac{24}{2+\varepsilon}.
\]
For sufficiently small nonzero $\varepsilon$, the quadratic roots
$b,c$ are real, distinct, greater than one, and close to $3,4$.
The sum remains $9$ and the product remains $24$.
Consequently $Z(0)$, $Z(-1)$, $Z'(0)$, the residue at one and the
zeta determinant are unchanged, while the full spectrum changes.
All eigenvalues outside this triple, including the zero eigenvalue,
remain fixed.  A determinant alone already permits the two-value
perturbation $2\mapsto2c$, $3\mapsto3/c$ for $c$ near $1$.
These countermodels concern the specified familiar invariants;
they make no claim about an arbitrarily encoded finite datum.
\end{proof}

\subsection{The fixed mixing recipe and its exact unknown-measure inverse}

\begin{theorem}[Existence and uniqueness of the linked mixing measure]
\label{final:O5}
Let $B\geq0$ be a self-adjoint operator on a Hilbert space.  The fixed
function $J(\lambda)=(1+\lambda)^{-2}$ always has the strong-operator
mixing representation
\begin{equation}\label{final:O5-J}
 \int_0^\infty e^{-sB}\,s e^{-s}\,ds
       =(\mathrm{Id}+B)^{-2}.
\end{equation}
If $B$ has a nonzero eigenvector for every $n\in\mathbb N_0$, then
the unknown finite complex Borel measure $\nu$ on $[0,\infty)$ in
\begin{equation}\label{final:O5-unknown}
 \int_{[0,\infty)}e^{-sB}\,\nu(ds)=(\mathrm{Id}+B)^{-2}
\end{equation}
is uniquely $\nu(ds)=s e^{-s}\,ds$.
Finite total variation is understood for a complex measure.  A single
eigenvector at every integer suffices; multiplicities and additional
spectrum do not affect the conclusion.

There is also a full-line formulation without initially defining
negative-time operators.  If $\nu$ is a finite positive Borel measure
on $\mathbb R$ and its finite scalar integer integrals satisfy
\begin{equation}\label{final:O5-integer-samples}
 \int_{\mathbb R}e^{-ns}\,\nu(ds)=(n+1)^{-2}
 \qquad(n\in\mathbb N_0),
\end{equation}
then $\nu(ds)=\mathbf1_{(0,\infty)}s e^{-s}\,ds$.
In this formulation mass, positive support and absence of an atom
at zero are conclusions of the supplied samples.
\end{theorem}

\begin{proof}
The scalar gamma integral gives, for every $\lambda\geq0$,
\[
 \int_0^\infty se^{-(1+\lambda)s}\,ds=(1+\lambda)^{-2}.
\]
The semigroup is a contraction, so each vector-valued integral is
absolutely integrable.  Integration against the spectral measure,
justified by bounded domination, proves
\eqref{final:O5-J} for every nonnegative self-adjoint $B$.

For uniqueness, apply \eqref{final:O5-unknown} to a nonzero
$n$-eigenvector.  This gives all scalar equations
\eqref{final:O5-integer-samples} on $[0,\infty)$.
Under $y=e^{-s}$, their left sides are the moments of the finite
complex pushforward measure on $(0,1]$, extended by zero to $[0,1]$.
The finite positive measure with density $-\log y$ on $(0,1)$ has
these same moments, since
\[
 \int_0^1y^n(-\log y)\,dy=(n+1)^{-2}.
\]
Uniform polynomial approximation on $[0,1]$ identifies the two
finite measures.  Undoing the measurable bijection $s=-\log y$
gives $\nu(ds)=se^{-s}\,ds$.  In particular the zeroth equation
recovers its mass; its lack of an atom at $s=0$ follows as well.

For a positive measure initially on all of $\mathbb R$, mass on
$(-\infty,-\varepsilon]$ would imply
\[
 (n+1)^{-2}\geq e^{n\varepsilon}
       \nu(( -\infty,-\varepsilon])
\]
for every $n$, which is impossible for positive mass.  Taking the
union over $\varepsilon=1/k$ proves support in $[0,\infty)$.
The preceding compact moment proof then applies.  No integral of
unbounded negative-time operators was assumed in this support step.
\end{proof}

\begin{theorem}[A Bernstein function is determined by every integer tail]
\label{final:O6}
Let $f,g$ be Bernstein functions, in their usual representation
category
\begin{equation}\label{final:O6-bernstein-representation}
 f(\lambda)=k_f+d_f\lambda+
   \int_{(0,\infty)}(1-e^{-\lambda s})\,\Pi_f(ds),
 \qquad \lambda\geq0,
\end{equation}
where $k_f,d_f\geq0$, $\Pi_f$ is positive, and
$\int(1\wedge s)\Pi_f(ds)<\infty$, and likewise for $g$.
For any fixed $N\in\mathbb N_0$, equality $f(n)=g(n)$ for every
integer $n\geq N$ implies $f=g$ on $[0,\infty)$, including equality
of killing, drift and L\'evy measure.
In particular the samples
\[
 f(n)=2\log(1+n)\qquad(n\geq N)
\]
identify the complete exponent $f(\lambda)=2\log(1+\lambda)$,
with $k_f=d_f=0$ and $\Pi_f(ds)=2e^{-s}ds/s$.
Knowledge of $f(A)$ on the marked integer eigenspaces supplies these
samples.  This inverse does not extend to arbitrary smooth functions.
\end{theorem}

The integer-tail uniqueness principle for Bernstein functions also appears
in Aguech and Jedidi, Corollary~4~\cite{AguechJedidi}.  The proof below is
independent and additionally records the recovery of the
L\'evy--Khintchine data in the present notation.

\begin{proof}
The finite positive measure
\[
 \rho_f=d_f\delta_1+
    (s\mapsto e^{-s})_\#[(1-e^{-s})\Pi_f(ds)]
 \quad\hbox{on }[0,1]
\]
has no atom at zero.  Its finiteness follows from
$1-e^{-s}\leq1\wedge s$.  Direct subtraction of
\eqref{final:O6-bernstein-representation} gives
\[
 f(n+1)-f(n)=\int_0^1y^n\rho_f(dy),\qquad n\geq0.
\]
Equality of the integer tails therefore identifies all moments of
$y^N\rho_f$ and $y^N\rho_g$.  Compact moment uniqueness gives
$y^N\rho_f=y^N\rho_g$.  Since the weight is strictly positive on
$(0,1]$, division on $[1/j,1]$ and passage to their increasing union
give equality there.  The absence of atoms at zero gives
$\rho_f=\rho_g$ on $[0,1]$.
The atom at one recovers the drift $d_f=d_g$.  The restriction to
$(0,1)$, the inverse coordinate $s=-\log y$, and division by the
strictly positive weight $1-e^{-s}$ recover
$\Pi_f=\Pi_g$.  With these fixed, a single common value at $N$
recovers $k_f=k_g$; this also works when $N>0$.

For the displayed target measure, the L\'evy integrability condition
holds at both endpoints.  Its exponent is zero at $\lambda=0$ and
has derivative
\[
 \frac{d}{d\lambda}\int_0^\infty
       (1-e^{-\lambda s})\frac{2e^{-s}}s\,ds
 =2\int_0^\infty e^{-(1+\lambda)s}\,ds
 =\frac2{1+\lambda}.
\]
Thus it equals $2\log(1+\lambda)$, proving the target assertion.
Functional calculus gives
$f(A)e_n=f(n)e_n$ and
\[
 D(f(A))=\left\{\sum c_ne_n:\sum|f(n)|^2|c_n|^2<\infty\right\}.
\]
Consequently equality on the marked integer eigenspaces has exactly
the required scalar meaning.  For arbitrary smooth functions the
nonzero perturbation $\varepsilon\sin^2(\pi\lambda)$ is invisible
at every nonnegative integer; adding it to the nonnegative target
also preserves nonnegativity.  It therefore gives a distinct scalar
function with the same functional calculus on $A$, outside the
Bernstein inverse category just proved.
\end{proof}

\begin{proposition}[Two different functional-calculus inverse questions]
\label{final:O6-functional-calculus}
For a fixed strictly increasing Bernstein function $f$, the full
self-adjoint operator $f(B)$ determines a nonnegative self-adjoint
$B$.  For the canonical $A$, the complete operator
$2\log(\mathrm{Id}+A)$ is equivalently determined by
$J(A)=(\mathrm{Id}+A)^{-2}$.
These statements retain the known scalar function $f$ or $J$.
\end{proposition}

\begin{proof}
Strict monotonicity and continuity give a Borel inverse on the range
of $f$.  The Borel functional calculus gives
$B=f^{-1}(f(B))$, with its exact spectral domain.  If an endpoint of
the range is a spectral limit, the spectral projection at that
unattained endpoint is zero, so the unbounded inverse is still
well-defined.  In the indicated special case, $J(A)$ has eigenvalues
$(1+n)^{-2}>0$, and
$-\log J(A)=2\log(\mathrm{Id}+A)$ on its functional-calculus domain.
Inverting a known function of an unknown operator is thus different
from Theorem~\ref{final:O6}, which identifies an unknown function
from its values on a known integer spectrum.
\end{proof}

\begin{theorem}[The Gamma shapes two and three retain all spectral and mixing data]
\label{final:O5-isospectral-boundary}
For $\alpha\in\{2,3\}$ set
\[
 p_\alpha(t)=\frac{t^{\alpha-1}e^{-t}}{\Gamma(\alpha)},
 \qquad \mathcal H_\alpha=L^2((0,\infty),p_\alpha(t)\,dt),
\]
and let $A_\alpha$ be the closure of
\[
 -tf''-(\alpha-t)f',\qquad D=C_c^\infty(0,\infty).
\]
Both endpoints are limit point in both cases.  The operators have
the same simple complete spectrum $\mathbb N_0$ and are intertwined
by a unitary carrying $1$ to $1$.  Therefore all their heat traces,
shifted zeta functions, resolvent trace functions and spectral
determinants agree, with the same conventions and domains.
Both satisfy \eqref{final:O5-unknown} with the same unique mixing
measure $\nu_2(ds)=se^{-s}\,ds$.  Their invariant coordinate laws
nevertheless differ: they are $p_2\,dt$ and
$p_3\,dt$, with means $2$ and $3$ respectively.
\end{theorem}

\begin{proof}
The divergence-form coefficients are $q_\alpha=tp_\alpha$.
Apart from an irrelevant nonzero constant, the second zero-energy
solution is $\int_1^t e^s s^{-\alpha}\,ds$.
At zero it is asymptotic to
$-t^{1-\alpha}/(\alpha-1)$; its weighted square is asymptotic to a
positive multiple of $t^{1-\alpha}$, nonintegrable for both
$\alpha=2$ and $3$.  At infinity the solution is asymptotic to
$e^t/t^\alpha$, whose weighted square is a positive multiple of
$e^t/t^{\alpha+1}$, also nonintegrable.
The constant solution is square integrable at both ends.
The maximal-domain and endpoint arguments of Theorem~\ref{final:O3}
therefore apply to each, giving their unique self-adjoint closures.
The form argument also applies, since the lower-cutoff energy is
$O(\varepsilon^{\alpha-1})$, proving conservativity and invariance
of $p_\alpha\,dt$.

The Rodrigues polynomials
\[
 L_n^{(\alpha-1)}(t)=
   \frac{t^{1-\alpha}e^t}{n!}
      \frac{d^n}{dt^n}(e^{-t}t^{n+\alpha-1})
\]
satisfy $A_\alpha L_n^{(\alpha-1)}=nL_n^{(\alpha-1)}$ because their
coefficients obey
$(k+1)(k+\alpha)a_{k+1}+(n-k)a_k=0$.
Integration by parts gives squared norm
$\Gamma(n+\alpha)/(n!\Gamma(\alpha))$ in $\mathcal H_\alpha$.
The weighted Laplace-transform proof of completeness applies without
change: $p_\alpha$ has the same exponential tail and all polynomial
moments.  Hence
\[
 e_{n,\alpha}=
 \sqrt{\frac{n!\Gamma(\alpha)}{\Gamma(n+\alpha)}}
       L_n^{(\alpha-1)}
\]
is a complete orthonormal eigenbasis, with $e_{0,\alpha}=1$.
The unitary $Ue_{n,2}=e_{n,3}$ intertwines the two operators and
fixes the distinguished constant vectors.  Every listed scalar
spectral invariant therefore agrees.  Theorem~\ref{final:O5}
gives the common mixing identity and its mixing-measure uniqueness.
Finally gamma integration gives
$\int t p_\alpha(t)\,dt=\alpha$, proving that the coordinate
probabilities differ despite all those agreements.
\end{proof}

\begin{corollary}[Canonical mixing construction is not unmarked coordinate recovery]
\label{final:O5-canonical-versus-identification}
The fixed exponent-two recipe $J$ constructs the Gamma$(2,1)$ mixing
measure for every nonnegative self-adjoint operator.  Integer
eigenvalues make that mixing measure uniquely recoverable from the
restricted operator identity.  They do not identify the stationary
coordinate law of an arbitrary realization, even when the constant
vector is retained.  To identify that candidate law one must retain
either its explicit equality to the unknown mixing measure in
\eqref{final:O5-unknown}, the local Pearson realization of
Theorem~\ref{final:O2}, or equivalent coordinate information.
\end{corollary}

\begin{proof}
Existence and mixing uniqueness are Theorem~\ref{final:O5}; the two
realizations in Theorem~\ref{final:O5-isospectral-boundary} are a
counterexample to the unmarked coordinate inverse, even under the
conjunction of all the stated spectral and mixing data.
For contrast, a fully marked package
$(\mathcal H,A,1,M_t)$ retains the multiplication operator
$M_tf(t)=tf(t)$ on $\{f:tf\in\mathcal H\}$, and recovers its
coordinate probability by
\[
 P(E)=\langle\mathbf1_E(M_t)1,1\rangle
\]
for every Borel $E$.  A unitary intertwining this multiplication
operator as well as fixing $1$ would preserve that probability,
so the preceding unitary cannot do so.
Furthermore, the different retained recipe
$J_a(\lambda)=(1+\lambda)^{-a}$, $a>0$, has mixing measure
$s^{a-1}e^{-s}\,ds/\Gamma(a)$ by the same gamma integral.
Spectral data do not choose this exponent.

In particular let $F$ forget a stationary coordinate realization
and retain its spectrum, and let $G_J$ construct the canonical
Gamma$(2,1)$ realization using the fixed recipe.  Then
\[
 G_J(F(p_3,A_3))=(p_2,A_2)\ne(p_3,A_3).
\]
Thus this construction is not a left inverse to forgetting on the
class of stationary coordinate realizations.  An identification
graph can use the linked unknown-measure theorem or retain the
fixed coordinate template in its nodes; a cycle through an
unmarked spectrum and canonical selection alone does not prove
identification of its original candidate.
\end{proof}
